\item {\bf Explainability of Logic-Based Systems}

The notion that AI systems should be explainable and their decisions interpretable has gained traction. 
For instance, \href{https://gdpr-text.com/read/article-15/#related_gdpr-a-15_1h}{GDPR’s Article 15} requires that individuals be provided 
\href{https://academic.oup.com/idpl/article/7/4/233/4762325}{“meaningful information”} about the logic of automated decisions. 
Independently of legal considerations, explainable AI protects the interests of both end users 
(the people employing AI to make or inform decisions) and to stakeholders (individuals affected by automated decisions). T
hese are some of the considerations cited in the literature:

\begin{enumerate}
  \item Respect: when decisions that affect you are made intelligible to you, this is considered an expression of respect for your 
  status as a human being whose interests are worthy of consideration.
  \item Assessing fairness of rules: knowing what rules are applied to reach a decision, we have the ability to evaluate whether the 
  rules are fair and whether they are fairly applied to us.
  \item Contesting and correcting decisions: understanding a decision we disagree with, allows us to contest it.
  \item Ability to change user behavior: understanding a decision that has a negative impact on us allows us to adapt our 
  behavior in order to obtain different results in the future.
\end{enumerate}


Explanability has many different meanings (see this paper by 
\href{https://ir.lawnet.fordham.edu/cgi/viewcontent.cgi?article=5569&context=flr}{Andrew Selbst and Solon Barocas}; and this paper by 
\href{https://arxiv.org/pdf/1702.08608.pdf}{Finale Doshi Velez and Been Kim}). Intuitively, logic-based systems should have more explanability power 
than machine learning-based “black-box” systems.  In this problem, let us explore what kind of explanability logic-based systems might 
provide and which human interests are protected.

Consider the Liar’s Puzzle from Problem 3.  Here, the first-order logic system takes in a set of facts and reasons over them to produce a conclusion. 
In this case, the conclusion judges the truthfulness of each person, which could carry with it real-world consequences. 
For example, someone might be fired if they are found guilty of crashing the server, whether they did it or not 
(AI is used widely for the allocation of healthcare resources, credit scoring, insurance, and content moderation, among other consequential domains).


\begin{enumerate}

  \input{06-explainability/01-question-1}
  \input{06-explainability/02-question-2}
  \input{06-explainability/03-question-3}
  
\end{enumerate}
